\chapter{Kernel Configuration}
\label{ch:kernelconfig}

The Linux kernel offers over 15,000 configuration options across dozens
of subsystems. Veilbox uses a carefully curated subset of approximately
111 options (plus the defconfig baseline) that provides exactly the
features needed for a virtualized container host.

\section{Configuration Strategy}

\subsection{Layered Configuration Model}

Veilbox uses a three-layer configuration model:

\begin{enumerate}[nosep]
  \item \textbf{Base layer}: \texttt{make defconfig} generates a
        configuration with approximately 10,000 options enabled. This
        provides a working baseline for x86\_64 hardware.
  \item \textbf{Veilbox layer}: The custom fragment
        (\texttt{custom-os.config}) adds or removes approximately 111
        options. This layer enables container-related features (namespaces,
        cgroups, overlay filesystem) and disables unnecessary hardware
        support.
  \item \textbf{Hardware layer}: Runtime-detected. The kernel uses
        module autoloading and udev (via devtmpfs) to load drivers for
        detected hardware. Since Veilbox uses an embedded initramfs with
        no module loading, all needed drivers must be compiled in (\texttt{y}).
\end{enumerate}

\subsection{Decision Rules}

Each configuration option in the Veilbox fragment follows these rules:

\begin{itemize}[nosep]
  \item \textbf{Essential features} (\texttt{y}): Options needed for system
        boot, container operation, or debugging are set to \texttt{y}
        (built-in).
  \item \textbf{Unnecessary features} (\texttt{n}): Hardware drivers for
        subsystems not present in a VM (sound, USB, Wi-Fi, Bluetooth,
        etc.) are explicitly disabled.
  \item \textbf{Modules} (\texttt{m}): Not used. Veilbox has no module
        loading infrastructure (no depmod, modprobe, or module directory).
        All drivers are built into the kernel.
\end{itemize}

\section{Critical Configuration Options}

\subsection{Container Runtime Support}

\begin{lstlisting}
CONFIG_NAMESPACES=y
CONFIG_UTS_NS=y
CONFIG_IPC_NS=y
CONFIG_PID_NS=y
CONFIG_NET_NS=y
CONFIG_USER_NS=y
CONFIG_CGROUPS=y
CONFIG_CGROUP_CPUACCT=y
CONFIG_CGROUP_DEVICE=y
CONFIG_CGROUP_FREEZER=y
CONFIG_CGROUP_SCHED=y
CONFIG_CPUSETS=y
CONFIG_MEMCG=y
\end{lstlisting}

These options enable Linux namespace isolation and cgroup resource
management. Namespaces provide process isolation (each container sees
its own PID tree, network stack, mount table, etc.). Cgroups provide
resource limits (CPU, memory, device access).

\subsection{Overlay Filesystem}

\begin{lstlisting}
CONFIG_OVERLAY_FS=y
\end{lstlisting}

The OverlayFS filesystem enables container image layering. Each container
image is composed of read-only layers with a thin writable layer on top.
This is the foundation of Docker/containerd image management.

\subsection{Networking}

\begin{lstlisting}
CONFIG_VETH=y
CONFIG_BRIDGE=y
CONFIG_NETFILTER=y
CONFIG_NETFILTER_XTABLES=y
CONFIG_IP_NF_FILTER=y
CONFIG_IP_NF_NAT=y
CONFIG_IP_NF_TARGET_MASQUERADE=y
CONFIG_VIRTIO_NET=y
\end{lstlisting}

Virtual Ethernet (veth) pairs connect container network namespaces to the
host bridge. Netfilter enables NAT and packet filtering for container
outbound traffic. The VirtIO net driver provides high-performance
paravirtualized networking in QEMU.

\subsection{Storage}

\begin{lstlisting}
CONFIG_VIRTIO_BLK=y
CONFIG_EXT4_FS=y
CONFIG_TMPFS=y
CONFIG_SQUASHFS=y
CONFIG_DEVTMPFS=y
\end{lstlisting}

VirtIO block provides paravirtualized storage in QEMU. ext4 is the primary
filesystem for the VEILBOX partition. tmpfs is used by the kernel for
the initramfs root. devtmpfs provides automatic device node creation.

\subsection{Serial Console}

\begin{lstlisting}
CONFIG_TTY=y
CONFIG_SERIAL_8250=y
CONFIG_SERIAL_8250_CONSOLE=y
CONFIG_SERIAL_CORE=y
CONFIG_SERIAL_CORE_CONSOLE=y
\end{lstlisting}

The serial console (ttyS0) is the primary management interface for headless
operation. The 8250 UART driver is the standard serial port controller in
x86\_64 systems.

\subsection{Mandatory for Boot}

\begin{lstlisting}
CONFIG_EFI_STUB=y
CONFIG_CMDLINE_BOOL=y
CONFIG_CMDLINE="console=ttyS0 console=tty1"
CONFIG_INITRAMFS_SOURCE="initramfs.list"
CONFIG_BLK_DEV_INITRD=y
\end{lstlisting}

\texttt{CONFIG\_EFI\_STUB} is required for QEMU's \texttt{-kernel} boot
mode — it produces a valid MZ/PE header that QEMU's firmware recognizes.
The kernel command line specifies both serial and VGA consoles.

\section{Configuration File Format}

The kernel configuration file (.config) uses a simple key-value format:

\begin{lstlisting}
# CONFIG_SOUND is not set
CONFIG_NAMESPACES=y
CONFIG_CGROUPS=y
CONFIG_EXT4_FS=m
\end{lstlisting}

The three value types are:
\begin{itemize}[nosep]
  \item \texttt{y}: Built into the kernel binary. Required for drivers
        that are needed at boot time (storage, console, initramfs).
  \item \texttt{m}: Built as a loadable kernel module. Not used in
        Veilbox.
  \item \texttt{n}: Disabled. The option and its code are excluded from
        the build entirely.
\end{itemize}

\section{Defconfig Reference}

The \texttt{make defconfig} target generates a configuration file that
enables the majority of options suitable for the target architecture.
On x86\_64, this includes:

\begin{itemize}[nosep]
  \item All x86\_64 CPU features: MMX, SSE, AVX, NX, PGE, MTRR, PAT,
        and so on.
  \item All major filesystems (ext4, XFS, btrfs, FAT, NTFS, etc.)
  \item Network protocol support (TCP, UDP, SCTP, DCCP, etc.)
  \item All common device drivers (storage, network, input, video, etc.)
  \item Security features (LSM, SELinux, AppArmor, IMA, EVM, etc.)
\end{itemize}

The Veilbox fragment reduces this to only the features needed for a
container host running in a virtual machine.

\section{Validating Configuration}

After applying the fragment and resolving dependencies, the configuration
should be validated:

\begin{lstlisting}[language=bash]
# Check for unresolved dependencies
make olddefconfig

# Check specifically for container support
grep -E "(NAMESPACES|CGROUPS|OVERLAY_FS)" .config

# Check for VirtIO drivers
grep -E "(VIRTIO_NET|VIRTIO_BLK)" .config

# Verify initramfs source path
grep CONFIG_INITRAMFS_SOURCE .config
\end{lstlisting}


\section{Detailed Option Reference}

\subsection{Namespace Configuration}

Linux namespaces are the foundation of container isolation. Each namespace
type isolates a different global resource:

\begin{longtable}{|p{3cm}|p{4cm}|p{5cm}|}
\hline
\textbf{Namespace} & \textbf{Isolates} & \textbf{Kernel Option} \\
\hline
Mount (mnt) & Filesystem mount points & CONFIG\_UTS\_NS \\
PID (pid) & Process IDs & CONFIG\_PID\_NS \\
Network (net) & Network stack & CONFIG\_NET\_NS \\
IPC (ipc) & System V IPC, POSIX messages & CONFIG\_IPC\_NS \\
UTS (uts) & Hostname, domain name & CONFIG\_UTS\_NS \\
User (user) & UID/GID mappings & CONFIG\_USER\_NS \\
Cgroup (cgroup) & Cgroup root directory & CONFIG\_CGROUPS \\
Time (time) & Clock sources & CONFIG\_TIME\_NS \\
\hline
\end{longtable}

\subsection{Control Group Controllers}

Cgroups limit and account for resource usage. Each controller regulates
a specific resource:

\begin{longtable}{|p{3cm}|p{4cm}|p{5cm}|}
\hline
\textbf{Controller} & \textbf{Resource} & \textbf{Kernel Option} \\
\hline
cpu & CPU scheduling & CONFIG\_CGROUP\_SCHED \\
cpuacct & CPU accounting & CONFIG\_CGROUP\_CPUACCT \\
cpuset & CPU/memory pinning & CONFIG\_CPUSETS \\
memory & Memory limits & CONFIG\_MEMCG \\
blkio & Block I/O limits & CONFIG\_BLK\_CGROUP \\
devices & Device access & CONFIG\_CGROUP\_DEVICE \\
freezer & Process freeze/thaw & CONFIG\_CGROUP\_FREEZER \\
pids & PID limits & CONFIG\_CGROUP\_PIDS \\
\hline
\end{longtable}

\subsection{Netfilter and NAT}

For container networking, Netfilter provides:

\begin{itemize}[nosep]
  \item \textbf{Packet filtering}: Accept, drop, or reject packets based
        on source/destination IP, port, protocol, and interface.
  \item \textbf{NAT (Network Address Translation)}: Masquerade container
        outbound traffic behind the host IP address.
  \item \textbf{Connection tracking}: Track TCP/UDP connection state for
        stateful filtering and NAT.
  \item \textbf{Forwarding}: Enable IP forwarding between the host
        network namespace and container namespaces.
\end{itemize}

The minimum Netfilter options for container networking:

\begin{lstlisting}
CONFIG_NETFILTER=y
CONFIG_NETFILTER_ADVANCED=y
CONFIG_NETFILTER_XTABLES=y
CONFIG_NETFILTER_XT_TARGET_MASQUERADE=y
CONFIG_NETFILTER_XT_MATCH_CONNTRACK=y
CONFIG_IP_NF_FILTER=y
CONFIG_IP_NF_NAT=y
CONFIG_IP_NF_TARGET_MASQUERADE=y
\end{lstlisting}
